% GENERATED FILE. Edit canonical lessons, then run npm run generate:books.
% canonical-source-hash: fnv1a64:9bde2d7316417567
% canonical-lessons: MR-C445-ankur, MR-C445-rujne, MR-C445-phulne, MR-C445-baharne, MR-C445-komejne

\chapter{Blossom and Fall}
\label{ch:mr-blossom-and-fall}
\hlchaptermodality{445}

\begin{chapteropening}
\textbf{By the end of this chapter:} \emph{I can talk about sprouts, and say plants take root, bloom, blossom and wither.}

Complete the last lesson of chapter 445: I can talk about sprouts, and say plants take root, bloom, blossom and wither.
\end{chapteropening}

\section[\texorpdfstring{\mr{अंकुर} (aṅkur)}{अंकुर (aṅkur)}]{\mr{अंकुर} (aṅkur) — a sprout; a shoot}

\label{lesson:MR-C445-ankur}



\begin{quote}
Before the new one: say the Marathi for to seep, then the Marathi for to overflow.
\end{quote}

\subsection*{You'll want to know: \mr{अंकुर}}
\textbf{\mr{अंकुर}} — \emph{aṅkur} — \textquotedblleft{}a sprout; a shoot\textquotedblright{}.

\begin{quote}
\mr{बीमधून} \mr{हिरवा} \mr{अंकुर} \mr{बाहेर} \mr{आला}.
\end{quote}

\begin{quote}
\emph{bīmadhūn hirvā aṅkur bāher ālā.}
\end{quote}

\begin{quote}
A green shoot came out of the seed.
\end{quote}

Masculine.

\mr{अंकुर} \mr{फुटणे}, a shoot comes up.

\begin{cousinweb}
\textbf{\mr{अंकुर}} (\emph{aṅkur}) is a tatsama, taken from Sanskrit \emph{aṅkura}, \textquotedblleft{}a sprout, a shoot\textquotedblright{} (Monier-Williams).
\end{cousinweb}

\begin{grammarlens}[title={What you've built}]
One more word for getting around.
\end{grammarlens}

\subsection*{Your turn}
\begin{itemize}
\raggedright
  \item \emph{Say it:} \emph{aṅkur}
  \item \emph{Say it:} \emph{aṅkur}, once more
  \item \emph{Say it:} \emph{bīmadhūn hirvā aṅkur bāher ālā.} — A green shoot came out of the seed.
  \item \emph{Recall:} say the Marathi for to seep, then the Marathi for to overflow, then say \emph{aṅkur} again
\end{itemize}

\begin{tcolorbox}[breakable,colback=teal!4,colframe=teal!35!black,title={Before you move on}]
What does \mr{अंकुर} mean? (\textquotedblleft{}A sprout; a shoot\textquotedblright{}.) Say it once more.
\end{tcolorbox}
\section[\texorpdfstring{\mr{रुजणे} (rujṇe)}{रुजणे (rujṇe)}]{\mr{रुजणे} (rujṇe) — to take root}

\label{lesson:MR-C445-rujne}



\begin{quote}
Before the new one: say the Marathi for to overflow, then the Marathi for a sprout.
\end{quote}

\subsection*{You'll want to know: \mr{रुजणे}}
\textbf{\mr{रुजणे}} — \emph{rujṇe} — \textquotedblleft{}to take root\textquotedblright{}.

\begin{quote}
\mr{पावसानंतर} \mr{रोपं} \mr{लवकर} \mr{रुजली}.
\end{quote}

\begin{quote}
\emph{pāvsānantar ropaṁ lavkar rujlī.}
\end{quote}

\begin{quote}
After the rain the saplings took root quickly.
\end{quote}

Past (he …): \mr{तो} \mr{रुजला}, \emph{to rujlā}.

Also for ideas and habits: \mr{सवय} \mr{रुजली}.

\begin{cousinweb}
\textbf{\mr{रुजणे}} (\emph{rujṇe}) is inherited: CDIAL lists Marathi \emph{ruj̈ṇẽ}, \textquotedblleft{}to sprout, grow\textquotedblright{}, under \emph{*ruhyati}, \textquotedblleft{}grows\textquotedblright{} (CDIAL 10797).
\end{cousinweb}

\begin{grammarlens}[title={What you've built}]
One more word for getting around.
\end{grammarlens}

\subsection*{Your turn}
\begin{itemize}
\raggedright
  \item \emph{Say it:} \emph{rujṇe}
  \item \emph{Say it:} \emph{rujṇe}, once more
  \item \emph{Say it:} \emph{pāvsānantar ropaṁ lavkar rujlī.} — After the rain the saplings took root quickly.
  \item \emph{Recall:} say the Marathi for to overflow, then the Marathi for a sprout, then say \emph{rujṇe} again
\end{itemize}

\begin{tcolorbox}[breakable,colback=teal!4,colframe=teal!35!black,title={Before you move on}]
What does \mr{रुजणे} mean? (\textquotedblleft{}To take root\textquotedblright{}.) Say it once more.
\end{tcolorbox}
\section[\texorpdfstring{\mr{फुलणे} (phulṇe)}{फुलणे (phulṇe)}]{\mr{फुलणे} (phulṇe) — to bloom}

\label{lesson:MR-C445-phulne}



\begin{quote}
Before the new one: say the Marathi for a sprout, then the Marathi for to take root.
\end{quote}

\subsection*{You'll want to know: \mr{फुलणे}}
\textbf{\mr{फुलणे}} — \emph{phulṇe} — \textquotedblleft{}to bloom\textquotedblright{}.

\begin{quote}
\mr{वसंतात} \mr{बागेत} \mr{गुलाब} \mr{फुलले}.
\end{quote}

\begin{quote}
\emph{vasantāt bāget gulāb phulle.}
\end{quote}

\begin{quote}
In spring the roses bloomed in the garden.
\end{quote}

Past (he …): \mr{तो} \mr{फुलला}, \emph{to phullā}.

Flowers \mr{फुलतात}; a happy face \mr{फुलतो} too.

\begin{cousinweb}
\textbf{\mr{फुलणे}} (\emph{phulṇe}) is inherited: CDIAL lists Marathi \emph{phulṇẽ}, \textquotedblleft{}to swell, blossom\textquotedblright{}, under Sanskrit \emph{phullati}, \textquotedblleft{}expands, opens\textquotedblright{} (CDIAL 9093).
\end{cousinweb}

\begin{grammarlens}[title={What you've built}]
One more word for getting around.
\end{grammarlens}

\subsection*{Your turn}
\begin{itemize}
\raggedright
  \item \emph{Say it:} \emph{phulṇe}
  \item \emph{Say it:} \emph{phulṇe}, once more
  \item \emph{Say it:} \emph{vasantāt bāget gulāb phulle.} — In spring the roses bloomed in the garden.
  \item \emph{Recall:} say the Marathi for a sprout, then the Marathi for to take root, then say \emph{phulṇe} again
\end{itemize}

\begin{tcolorbox}[breakable,colback=teal!4,colframe=teal!35!black,title={Before you move on}]
What does \mr{फुलणे} mean? (\textquotedblleft{}To bloom\textquotedblright{}.) Say it once more.
\end{tcolorbox}
\section[\texorpdfstring{\mr{बहरणे} (baharṇe)}{बहरणे (baharṇe)}]{\mr{बहरणे} (baharṇe) — to blossom}

\label{lesson:MR-C445-baharne}



\begin{quote}
Before the new one: say the Marathi for to take root, then the Marathi for to bloom.
\end{quote}

\subsection*{You'll want to know: \mr{बहरणे}}
\textbf{\mr{बहरणे}} — \emph{baharṇe} — \textquotedblleft{}to blossom\textquotedblright{}.

\begin{quote}
\mr{आंब्याचं} \mr{झाड} \mr{फुलांनी} \mr{बहरलं}.
\end{quote}

\begin{quote}
\emph{āmbyācaṁ jhāḍ phulānnī baharlaṁ.}
\end{quote}

\begin{quote}
The mango tree blossomed all over.
\end{quote}

Past (he …): \mr{तो} \mr{बहरला}, \emph{to baharlā}.

A whole tree in flower \mr{बहरतो}; the noun is \mr{बहर}.

\begin{cousinweb}
\textbf{\mr{बहरणे}} (\emph{baharṇe}) is built on \textbf{\mr{बहर}}, \textquotedblleft{}full bloom\textquotedblright{}, a loan usually traced to Persian \emph{bahār}, \textquotedblleft{}spring\textquotedblright{}.
\end{cousinweb}

\begin{grammarlens}[title={What you've built}]
One more word for getting around.
\end{grammarlens}

\subsection*{Your turn}
\begin{itemize}
\raggedright
  \item \emph{Say it:} \emph{baharṇe}
  \item \emph{Say it:} \emph{baharṇe}, once more
  \item \emph{Say it:} \emph{āmbyācaṁ jhāḍ phulānnī baharlaṁ.} — The mango tree blossomed all over.
  \item \emph{Recall:} say the Marathi for to take root, then the Marathi for to bloom, then say \emph{baharṇe} again
\end{itemize}

\begin{tcolorbox}[breakable,colback=teal!4,colframe=teal!35!black,title={Before you move on}]
What does \mr{बहरणे} mean? (\textquotedblleft{}To blossom\textquotedblright{}.) Say it once more.
\end{tcolorbox}
\section[\texorpdfstring{\mr{कोमेजणे} (komejṇe)}{कोमेजणे (komejṇe)}]{\mr{कोमेजणे} (komejṇe) — to wither}

\label{lesson:MR-C445-komejne}



\begin{quote}
Before the new one: say the Marathi for to bloom, then the Marathi for to blossom.
\end{quote}

\subsection*{You'll want to know: \mr{कोमेजणे}}
\textbf{\mr{कोमेजणे}} — \emph{komejṇe} — \textquotedblleft{}to wither\textquotedblright{}.

\begin{quote}
\mr{पाणी} \mr{न} \mr{मिळाल्यामुळे} \mr{फुलं} \mr{कोमेजली}.
\end{quote}

\begin{quote}
\emph{pāṇī na miḷālyāmuḷe phulaṁ komejlī.}
\end{quote}

\begin{quote}
The flowers withered because they got no water.
\end{quote}

Past (he …): \mr{तो} \mr{कोमेजला}, \emph{to komejlā}.

Flowers and faces \mr{कोमेजतात}.

\begin{cousinweb}
\textbf{\mr{कोमेजणे}} (\emph{komejṇe}) is inherited: CDIAL lists Marathi \emph{komej̈ṇẽ}, \textquotedblleft{}wither\textquotedblright{}, under \emph{*kōmh} (CDIAL 3524).
\end{cousinweb}

\begin{grammarlens}[title={What you've built}]
One more word for getting around.
\end{grammarlens}

\subsection*{Your turn}
\begin{itemize}
\raggedright
  \item \emph{Say it:} \emph{komejṇe}
  \item \emph{Say it:} \emph{komejṇe}, once more
  \item \emph{Say it:} \emph{pāṇī na miḷālyāmuḷe phulaṁ komejlī.} — The flowers withered because they got no water.
  \item \emph{Recall:} say the Marathi for to bloom, then the Marathi for to blossom, then say \emph{komejṇe} again
\end{itemize}

\begin{tcolorbox}[breakable,colback=teal!4,colframe=teal!35!black,title={Before you move on}]
What does \mr{कोमेजणे} mean? (\textquotedblleft{}To wither\textquotedblright{}.) Say it once more.
\end{tcolorbox}
